\chapter*{引言}
\markboth{引言}{}
\addcontentsline{toc}{chapter}{引言}

这一系列笔记的目的是解释与几何朗兰兹对应有关的当前想法。该对应应当把范畴 $D(\Bun_G)$ 与叠 $\LocSys_{\Ch G}$ 上的拟凝聚层范畴联系起来。

本册研究这些猜想性构造所涉及的几何，更确切地说，所涉及的范畴。它们大多是经典自守形式论中某些函数空间的几何类比，例如主系列与 Whittaker 模型。

关于术语的一点说明：我们说「范畴」，指的是 [GL:DG] 所采用约定下的 DG 范畴；极限与余极限都在该文记号下的 $(\infty,1)$-范畴 $\mathrm{DGCat}_{cont}$ 中取。

关于 $\Bun_G$ 或 $\Bun_B$ 这类叠上 D-模范畴的定义，见 [DrGa1]。

\begin{annote}
本册在系列中的位置：§4 的谱侧讨论遵循 [GL:Eis]；§7--§8 的结论供 [GL:funct] 使用，后者据以对 $G=GL_2$ 构造几何朗兰兹函子 $\Psi_G$；§1 的拟定理 1.3.1 与 §5 的拟定理 5.3.3 分别在 [GL:contr]、[Ga2] 中证明。读者若只关心主角「扩展 Whittaker 范畴」，可先读 §6--§7，再回头补 §1--§4 的记号。\par\vspace{0.15em}{\raggedleft 【来源：原书各处引用（A）】\par}
\end{annote}
